\documentclass[10pt,a4paper]{amsart}
\usepackage[margin=19mm]{geometry}
\usepackage{amsmath,amssymb,mathtools}
\usepackage[scr=rsfs,cal=euler]{mathalfa}
\usepackage{xfrac}
\usepackage[unicode,hidelinks,bookmarks=false]{hyperref}
\newcommand{\ie}{i.e.\ }
\newcommand{\cf}{cf.\ }
\newcommand{\resp}{respectively\ }
\newcommand{\Subsection}{\subsection}
\providecommand{\nobreakdash}{\nobreak\hbox{-}}
\setlength{\emergencystretch}{2em}
\multlinegap0pt
\usepackage{bm}
\usepackage{bbm}
\usepackage{dsfont}
\usepackage{xspace}
\usepackage{cancel}
\usepackage{mathtools}
\usepackage{amsthm}
\makeatletter
\@ifundefined{theorem}{\newtheorem{theorem}{Theorem}}{}
\makeatother
\title[First order joint differential projective invariants]{First order joint differential projective invariants}
\author{Leonid Bedratyuk}
\date{Source: arXiv:2507.13051v2 (CC BY 4.0)}
\begin{document}
\maketitle

\noindent\emph{Source and licence.} First order joint differential projective invariants. Version arXiv:2507.13051v2 (CC BY 4.0), distributed under the Creative Commons Attribution 4.0 International licence, as recorded in the arXiv licence metadata for this exact version and shown on the abstract page https://arxiv.org/abs/2507.13051v2. Full rights notice: licenses/arxiv-2507.13051-CC-BY-4.0.txt.\par\smallskip

\noindent\textit{Editorial scope.} Counted unit: the Theorem whose source label is \texttt{None} in the article (source0.tex lines 1248--1258, source label \texttt{None}), delivered together with the article's own complete original proof of it (source0.tex lines 1260--1373, one \texttt{proof} environment of 114 lines). No mathematics is written, repaired or re-derived: statement and proof are the article's own text line for line. Required context retained uncounted: thm:first-prol (lines 279--296); thm:zn-general (lines 1183--1205). Every definition, display and label of those lines is retained unchanged. Omitted from this package and not redistributed: 1--278, 297--1182, 1206--1247, 1259--1259, 1374--1437. Nothing inside a retained excerpt is omitted or altered: every retained body is delivered exactly as the article's lines are written, so no omission receipt of any kind accompanies this package. Citations: the retained excerpts carry no citation command at all, so no bibliography entry of the article is transcribed and no reference is redistributed. Reference adaptation: none was needed and none was made; every reference inside the delivered unit resolves inside it, so no unresolved reference or citation is produced. Printed slips kept literal: no printed word, symbol, hypothesis or number of the counted statement or of its proof was altered. Layout and class: the article's own class is \texttt{amsart} with packages including latexsym, amsmath, amscd, amssymb, amsthm, longtable, fontenc, inputenc, babel, graphicx, xcolor, enumitem, blkarray; the wrapper is the lane's pinned \texttt{amsart} editorial wrapper at 10pt on A4, which restates the theorem-like environments the retained text needs and adds the article's own macro definitions that the retained text needs (none), with extra wrapper packages (bm, bbm, dsfont, xspace, cancel, mathtools). The delivered numbering is the unit's own. Assets: no retained span contains a figure environment, a graphics-inclusion command, a PDF-page inclusion, a table or a TikZ picture; no image file is copied into this package, the package declares no asset dependency and no image file is redistributed with it. Licence: version arXiv:2507.13051v2 of this article is distributed under the Creative Commons Attribution 4.0 International licence, as recorded in the arXiv licence metadata for this exact version and shown on the abstract page https://arxiv.org/abs/2507.13051v2. A full rights notice is delivered at licenses/arxiv-2507.13051-CC-BY-4.0.txt. Characters: every retained excerpt is pure ASCII, so the delivered engine, XeLaTeX with its default Latin Modern text font, reports no missing character.
\begin{theorem}\label{thm:first-prol}
Let $T$ be the projective transformation. Then the partial derivatives of the transformed function $\tilde u$ are given by:
\begin{align*}
\tilde u_{1,0}(\tilde x, \tilde y) &=
\frac{c_1 x + c_2 y + c_3}{D} \left(
  -\det\begin{vmatrix} b_1 & b_2 \\ c_1 & c_2 \end{vmatrix} (u_{1,0} x + u_{0,1} y)
  +\det\begin{vmatrix} b_2 & b_3 \\ c_2 & c_3 \end{vmatrix} u_{1,0}
  +\det\begin{vmatrix} b_1 & b_3 \\ c_1 & c_3 \end{vmatrix} u_{0,1}
\right), %\label{eq:u10-tilde}
\\[4pt]
\tilde u_{0,1}(\tilde x, \tilde y) &=
\frac{c_1 x + c_2 y + c_3}{D} \left(
  \det\begin{vmatrix} a_1 & a_2 \\ c_1 & c_2 \end{vmatrix} (u_{1,0} x + u_{0,1} y)
  -\det\begin{vmatrix} a_2 & a_3 \\ c_2 & c_3 \end{vmatrix} u_{1,0}
  +\det\begin{vmatrix} a_1 & a_3 \\ c_1 & c_3 \end{vmatrix} u_{0,1}
\right).% \label{eq:u01-tilde}
\end{align*}
\end{theorem}

\begin{theorem}
\label{thm:zn-general}
Let $n\ge5$ and let  $g=\gcd(n,3)$.  
Define
$$
z_n\;=\begin{cases}
\frac{\displaystyle
      \Bigl(\prod_{i=5}^{\,n} \Delta_{12i}\Bigr)^{3}\;
      \bigl(\Delta_{134}\,\Delta_{234}\bigr)^{\,n-3}}
     {\displaystyle
      \bigl(\Delta_{123}\,\Delta_{124}\bigr)^{\,2n-9}}, \text{ if  } g=1,\\
			\frac{\displaystyle
      \Bigl(\prod_{i=5}^{\,n} \Delta_{12i}\Bigr)\;
      \bigl(\Delta_{134}\,\Delta_{234}\bigr)^{\,\frac{n}{3}-1}}
     {\displaystyle
      \bigl(\Delta_{123}\,\Delta_{124}\bigr)^{\,\frac{2n}{3}-3}},   \text{  if  } g=3\\
			\end{cases}
$$
Then \( z_n \) is a relative projective invariant, and its weight is equal
$$
\operatorname{wt}(z_n)=-\frac{1}{g}.
$$
\end{theorem}

\begin{theorem}
Let
\[
W_n := \left\{ \omega \in \mathbb{Q} \,\middle|\, \exists\, F \in \mathcal{I}_n \text{ with } \operatorname{wt}(F) = \omega \right\}
\]
be the additive group of all possible weights of first-order relative projective invariants for \( n \) points, and let \( g = \gcd(n, 3) \). Then
\[
W_n = \tfrac{1}{g} \mathbb{Z},
\]
and hence \( \mathcal{I}_n = \mathcal{I}_{n,0}(z_n) \).
\end{theorem}

\begin{proof}

Consider the one-parameter subgroup of \( PGL(3, \mathbb{R}) \) acting by
\[
(x_i, y_i) \mapsto \left( \frac{x_i}{1 - tx_i}, \frac{y_i}{1 - tx_i} \right).
\]
The Jacobian of this transformation is
\[
(1 - tx_i)^{-3}.
\]


Since \(F\in\mathcal{I}_n\) is, in general, a rational function in the variables
\(x_i,y_i,p_i,q_i\) $i=1,\ldots, n$, write \(F=P/Q\) with coprime polynomials \(P,Q\).
Fix an index \(i\) and consider the substitution induced by the prolonged action 
\[
(x_i, y_i, p_i, q_i) \mapsto \left(
\frac{x_i}{1 - tx_i},\;
\frac{y_i}{1 - tx_i},\;
-(1 - tx_i)\left( t p_i x_i + t q_i y_i - p_i \right),\;
(1 - tx_i) q_i
\right),
\]
from Theorem~\ref{thm:first-prol}.
By factoring out the maximal power of \((1-tx_i)\) common to all terms, there exist integers
\(m_P(i),m_Q(i)\in\mathbb{Z}\) and polynomials \(\widehat P_i,\widehat Q_i\) not divisible by \((1-tx_i)\) such that
\[
P \mapsto(1-tx_i)^{m_P(i)}\widehat P_i,\qquad
Q\mapsto(1-tx_i)^{m_Q(i)}\widehat Q_i.
\]
Hence
\[
F \mapsto(1-tx_i)^{m_P(i)-m_Q(i)}\cdot \widehat F_i,\qquad
\widehat F_i:=\widehat P_i/\widehat Q_i,
\]
and therefore the exponent \(m_i:=m_P(i)-m_Q(i)\) is an integer.
On the other hand, by the definition of weight and the Jacobian factor \((1-tx_i)^{-3}\),
the relative invariant \(F\) transforms as
\(
F \mapsto (1-tx_i)^{-3\omega}\cdot F.
\)
Comparing the exponents of the factor \((1-tx_i)\) for this fixed \(i\) yields
\[
-3\omega=m_i\in\mathbb{Z}.
\]
Since the index \(i\) was arbitrary, we conclude that \(-3\omega\in\mathbb{Z}\).


Now consider another one-parameter subgroup given by
\[
(x_i, y_i) \mapsto (e^t x_i, y_i).
\]
The Jacobian of this transformation is \( e^t \), so the action on a relative invariant \( F \) of weight \( \omega \) is
\[
F \mapsto e^{nt\omega} \cdot F.
\]
According to Theorem~\ref{thm:first-prol}, the prolonged action of this transformation is
\[
(x_i, y_i, p_i, q_i) \mapsto \left(e^t x_i,\; y_i,\; \frac{p_i}{e^t},\; q_i\right).
\]


Factoring out the maximal power of \(e^t\) common to all terms, there exist integers \(m_P(i),m_Q(i)\in\mathbb{Z}\) and polynomials \(\widehat P_i,\widehat Q_i\) (not divisible by \(e^t\)) such that
\[
P \mapsto e^{t\,m_P(i)}\,\widehat P_i,\qquad
Q \mapsto e^{t\,m_Q(i)}\,\widehat Q_i.
\]
Therefore
\[
F \mapsto e^{t\,\bigl(m_P(i)-m_Q(i)\bigr)}\cdot \widehat F_i,\qquad
\widehat F_i:=\widehat P_i/\widehat Q_i,
\]
and the exponent \(m_i:=m_P(i)-m_Q(i)\) is an integer. Comparing this with the weight transformation
\(
F \mapsto e^{n t \omega}\cdot F
\)
yields
\[
n\omega = m_i \in \mathbb{Z}.
\]
Since \(i\) was arbitrary, we conclude that \(n\omega\in\mathbb{Z}\).



Combining these results, we find that
\[
\omega \in \tfrac{1}{3} \mathbb{Z} \cap \tfrac{1}{n} \mathbb{Z} = \tfrac{1}{g} \mathbb{Z}, \quad \text{where } g= \gcd(n, 3),
\]
so
\[
W_n \subseteq \tfrac{1}{g} \mathbb{Z}.
\]

Theorem~\ref{thm:zn-general} guarantees the existence of a relative invariant \( z_n \in \mathcal{I}_n \) of weight \( -\tfrac{1}{g} \), which implies
\[
-\tfrac{1}{g} \in W_n \quad \Rightarrow \quad \tfrac{1}{g} \mathbb{Z} \subseteq W_n.
\]
Thus, we obtain the equality:
\[
W_n = \tfrac{1}{g} \mathbb{Z}.
\]

Now consider an arbitrary relative invariant \( F \in \mathcal{I}_n \). From the above, its weight \( \operatorname{wt}(F) \) must be divisible by \( -\tfrac{1}{g} \), i.e., \( -g \cdot \operatorname{wt}(F) \in \mathbb{Z} \). Since \( z_n \) has weight \( -\tfrac{1}{g} \), the product
\[
F \cdot z_n^{-g \cdot \operatorname{wt}(F)}
\]
has weight zero, and hence belongs to the field of absolute invariants \( \mathcal{I}_{n,0} \). Therefore, we can write
\[
F = z_n^{-g \cdot \operatorname{wt}(F)} \cdot F'
\quad \text{for some } F' \in \mathcal{I}_{n,0}.
\]

This completes the proof that \( \mathcal{I}_n = \mathcal{I}_{n,0}(z_n) \).
\end{proof}

\end{document}
